\section{复合、隐式与反函数求导}
\label{sec:02-Calculus/03-Differentiation/03}

在一台两级传动仪器中，若拧动旋钮输入微小角度 \(h\)，中间轴按 \(3\) 倍比例转动，指针再按 \(0.5\) 倍比例响应中间轴。最终指针的偏移量为 \(1.5h\)。

这种``各级局部放大率相乘''的直观几何关系，正是\textbf{链式法则（Chain Rule）}的核心。链式法则不仅是复合函数求导的基本定理，更是深度学习中反向传播（Backpropagation）算法的数学基石。

\subsection{放大率相乘与链式法则}

设中间变量为 \(u=f(x)\)，外层变量为 \(y=g(u)\)。利用上一节的线性近似表示：
当输入产生微小扰动 \(h\) 时，中间层的响应增量为
\[
\Delta u = f'(x)h + o(h).
\]
将 \(\Delta u\) 作为输入代入外层函数 \(g\) 的线性模型中，可得总体增量为
\[
\Delta y = g'(u)\Delta u + o(\Delta u) = g'\bigl(f(x)\bigr)f'(x)h + o(h).
\]
保留 \(h\) 的一阶项系数，即可得到复合函数的求导公式：
\[
(g \circ f)'(x) = g'\bigl(f(x)\bigr)f'(x),
\qquad\text{莱布尼茨记号下即为}\qquad
\frac{dy}{dx} = \frac{dy}{du} \cdot \frac{du}{dx}.
\]

莱布尼茨记号展现了形式上的``分数约分''美感，其本质是两级局部伸缩率的累乘。

\begin{center}
\begin{tikzpicture}[x=1cm,y=1cm,font=\small,>=stealth]
  \draw[black!55] (0,2.4) -- (5.5,2.4);
  \draw[very thick,blue!70!black] (1.5,2.4) -- (2.0,2.4);
  \node[left] at (0,2.4) {输入层 \(x\)};
  \node[above,font=\footnotesize] at (1.75,2.45) {\(h\)};

  \draw[black!55] (0,1.2) -- (5.5,1.2);
  \draw[very thick,blue!70!black] (1.5,1.2) -- (3.0,1.2);
  \node[left] at (0,1.2) {中间层 \(u=f(x)\)};
  \node[above,font=\footnotesize] at (2.25,1.25) {\(\Delta u \approx f'(x)h\)};

  \draw[black!55] (0,0) -- (5.5,0);
  \draw[very thick,blue!70!black] (1.5,0) -- (2.25,0);
  \node[left] at (0,0) {输出层 \(y=g(u)\)};
  \node[above,font=\footnotesize] at (1.875,0.05) {\(\Delta y \approx g'(u)f'(x)h\)};

  \draw[->,black!70] (4.2,2.3) -- (4.2,1.35);
  \node[right,font=\footnotesize] at (4.3,1.82) {\(f\)：放大率 \(f'(x)\)};
  \draw[->,black!70] (4.2,1.1) -- (4.2,0.15);
  \node[right,font=\footnotesize] at (4.3,0.62) {\(g\)：放大率 \(g'(u)\)};
\end{tikzpicture}

\emph{扰动信号依次穿过各层，每层按当前工作点上的导数（局部放大率）进行伸缩。多层复合的总放大率即为各层放大率之积。}
\end{center}

例如求 \(y=\sin(x^2)\) 的导数：外层 \(\sin(u)\) 在 \(u=x^2\) 处的局部放大率为 \(\cos(x^2)\)，内层 \(u=x^2\) 的局部放大率为 \(2x\)，二者相乘即得 \(y'=2x\cos(x^2)\)。当结构扩充至多层（如 \(e^{-\sqrt{1+x^2}}\)）时，只需将信息流沿传递路径上的每一层局部导数逐项连乘即可。

\subsection{梯度消失与梯度爆炸的数学根源}

``放大率相乘''带来了一个极其重要的临界效应：\textbf{整个链条的敏感度取决于所有环节的共同作用；只要其中任意一环的局部导数趋于零，整条链条的信息传递就会中断。}

以药理学模型为例：设给药剂量为 \(x\)，血药浓度为 \(u=f(x)\)，疗效为 \(y=g(u)\)。若增加剂量能有效提升浓度（\(f'(x)\) 较大），但血药浓度已进入受体饱和的平台期（\(g'(u)\approx 0\)），则最终的边际疗效 \(\frac{dy}{dx}=g'(u)f'(x)\approx 0\)。

在深度神经网络中，这一现象体现得更为剧烈。一个 \(L\) 层深层网络对早期参数的梯度，包含了 \(L\) 个局部 Jacobi 矩阵一阶导数的连乘。若各层的平均伸缩因子小于 \(1\)，乘积将随层数呈指数级衰减，导致\textbf{梯度消失（Vanishing Gradient）}，前层参数无法更新；反之若平均伸缩因子大于 \(1\)，则会导致\textbf{梯度爆炸（Exploding Gradient）}。这并非工程实现上的偶然缺陷，而是链式法则连乘结构在多层非线性系统中的必然数学结果。

需要强调的是，链式法则成立的充要条件是：\textbf{内层函数在 \(x\) 处可导，且外层函数在对应的 \(u=f(x)\) 处可导}。以 \(h(x)=|x^2-1|\) 为例，虽然内层 \(u=x^2-1\) 处处可导，外层绝对值函数仅在原点不可导；但当 \(x=\pm 1\) 时，内层输出恰为 \(u=0\)，正好触碰外层的不可导尖角，导致 \(h(x)\) 在 \(x=\pm 1\) 处不可导。

\subsection{隐式约束与相关变化率}

对于方程 \(x^2+y^2=25\)，若要在点 \((3,4)\) 处求切线斜率，无需解出显式表达式 \(y=\sqrt{25-x^2}\)。

由于动点在曲线上移动时必须严格满足方程约束，我们可以直接将 \(y\) 视为 \(x\) 的隐式复合函数，对方程两边同时关于 \(x\) 求导：
\[
2x+2y\cdot\frac{dy}{dx}=0 \implies \frac{dy}{dx}=-\frac{x}{y}
\]
代入点 \((3,4)\)，即可得到切线斜率为 \(-\frac{3}{4}\)。其中对 \(y^2\) 项求导得到 \(2y\cdot\frac{dy}{dx}\)，正是将 \(y^2\) 视作外层平方、内层为 \(y(x)\) 的链式法则体现。

\begin{center}
\begin{tikzpicture}[x=0.55cm,y=0.55cm,font=\small,>=stealth]
  \draw[->,black!55] (-6.2,0) -- (7.2,0) node[right] {$x$};
  \draw[->,black!55] (0,-6.2) -- (0,6.2) node[above] {$y$};
  \draw[blue!70!black,thick] (0,0) circle (5);
  
  \draw[black,thick] (0.5,5.875) -- (5.5,2.125);
  \fill (3,4) circle (2.2pt);
  \node[above right] at (3,4) {$(3,4)$};
  \draw[->,red!80!black,thick] (3,4) -- (4.2,3.1);
  \node[above right,font=\footnotesize,red!80!black] at (3.8,3.4) {切向量};

  \draw[dashed,red!70!black,thick] (5,-5.2) -- (5,5.2);
  \fill (5,0) circle (2.2pt);
  \node[below right,font=\footnotesize] at (5,0) {$(5,0)$};
  \node[right,font=\footnotesize,red!70!black] at (5.2,2.5) {$y=0$：切线竖直};
\end{tikzpicture}

\emph{隐式方程将变量约束在特定的曲线上。在点 $(3,4)$ 处，局部允许的位移方向切线斜率为 $-x/y$；而在 $(5,0)$ 处，切线趋于竖直，分母为零表明此处无法将 $y$ 表示为 $x$ 的单值可微函数。}
\end{center}

式中分母在 $y=0$ 处归零（如点 $(5,0)$），表明该点附近的曲线切线为竖直状态，无法在此局部建立 $y=f(x)$ 的可微函数关系。这一几何直观在严谨分析中由\textbf{隐式函数定理（Implicit Function Theorem）}所保障，并进一步推广为多变量微积分中的约束曲面与 Lagrange 乘子法。

将隐式求导应用到时间自变量 $t$ 上，即可处理\textbf{相关变化率（Related Rates）}问题——即多个由某种物理或几何几何关系“锁死”的变量，当其中一个量随时间变化时，推算其他量随时间跟着变化的速率。

例如，一根长 \(5\,\mathrm{m}\) 的梯子靠在墙上，底端离墙距离为 \(x(t)\)，顶端离地高度为 \(y(t)\)。无论梯子如何滑动，两者必须时刻满足勾股定理约束：
\[
x(t)^2 + y(t)^2 = 25
\]
由于 \(x\) 与 \(y\) 均为时间 \(t\) 的函数，将约束方程两边对时间 \(t\) 应用链式法则求导：
\[
\frac{d}{dt}\left(x^2\right) + \frac{d}{dt}\left(y^2\right) = \frac{d}{dt}(25) 
\implies 2x\frac{dx}{dt} + 2y\frac{dy}{dt} = 0 
\implies x\dot{x} + y\dot{y} = 0
\]
此式即为两端运动速度（变化率 \(\dot{x}\) 与 \(\dot{y}\)）之间必须满足的约束关系。

若底端以 \(\dot{x} = 2\,\mathrm{m/s}\) 的速度向外拉动，当底端到达 \(x = 3\,\mathrm{m}\) 时，由原方程解得此时高度 \(y = \sqrt{25 - 3^2} = 4\,\mathrm{m}\)。将这一瞬间的具体数值代入速度约束式：
\[
3 \times 2 + 4 \times \dot{y} = 0 \implies \dot{y} = -1.5\,\mathrm{m/s}
\]
结果中的负号表示顶端高度 \(y\) 正在减少，即以 \(1.5\,\mathrm{m/s}\) 的速度沿墙面下滑。

在推导过程中，核心在于\textbf{必须先对动态过程求导建立变化率之间的关系，再代入特定瞬间的数值}；若先代入常数，方程将退化为静态恒等式 \(3^2+4^2=25\)，变化率信息便会丢失。

\subsection{反函数导数与斜率倒数关系}

设 \(y=f(x)\) 在局部单调可逆，其反函数为 \(x=f^{-1}(y)\)。

\subsubsection{1. 直观直觉：放大率与倒数关系}
导数的物理本质是\textbf{局部放大率}：
\begin{itemize}\tightlist
  \item 原函数 \(y=f(x)\) 表示：当 \(x\) 变化时，\(y\) 的变化率（放大倍数）为 \(\frac{dy}{dx}\)。
  \item 反函数 \(x=f^{-1}(y)\) 则反过来表示：当 \(y\) 变化时，\(x\) 的变化率（缩放倍数）为 \(\frac{dx}{dy}\)。
\end{itemize}
如果正向映射将输入放大 \(m\) 倍，为了在复合后恢复原状，逆映射必然将其缩放为原来的 \(\frac{1}{m}\) 倍。用莱布尼茨记号表达极为直观：
\[
\frac{dx}{dy} = \frac{1}{\frac{dy}{dx}}
\]

利用链式法则进行代数推导，只需对恒等式 \(f^{-1}\bigl(f(x)\bigr)=x\) 两边关于 \(x\) 求导：
\[
(f^{-1})'\bigl(f(x)\bigr)\cdot f'(x) = 1
\implies
(f^{-1})'(y) = \frac{1}{f'(x)} = \frac{1}{f'\bigl(f^{-1}(y)\bigr)}
\]

\subsubsection{2. 几何视角：镜像对称与斜率变换}
从几何上看，函数 \(y=f(x)\) 与其反函数 \(x=f^{-1}(y)\) 的图像关于直线 \(y=x\) 轴对称：
\begin{itemize}\tightlist
  \item 对称变换将原曲线上点 \((x_0, y_0)\) 的切线（斜率为 \(m\)），映射为反曲线上点 \((y_0, x_0)\) 的切线（斜率为 \(\frac{1}{m}\)）。
  \item \textbf{陡峭变平缓}：若原函数在某处切线极为陡峭（\(m > 1\)），翻转后的反函数在该点切线就会变得极其平缓（\(\frac{1}{m} < 1\)）。
\end{itemize}

\subsubsection{3. 可导条件与退化情形（\(f'(x) \neq 0\)）}
公式成立的前提条件是 \(f'(x) \neq 0\)，我们可以通过两类典型函数来理解：

\begin{itemize}\tightlist
  \item \textbf{正常情形（\(y=\mathrm{e}^{x}\) 与 \(x=\ln y\)）}：\\
  在 \(x=2\) 处，\(f'(2)=\mathrm{e}^{2} \approx 7.39\)（指数函数增长极快）；\\
  其反函数 \(x=\ln y\) 在 \(y=\mathrm{e}^{2}\) 处的导数为 \((\ln y)' = \frac{1}{\mathrm{e}^{2}} \approx 0.135\)（对数函数增长极慢）。

  \item \textbf{奇异退化（\(y=x^3\) 与 \(x=y^{1/3}\)）}：\\
  原函数 \(y=x^3\) 在 \(x=0\) 处满足 \(f'(0)=0\)，图像在原点处被``压平''（切线呈水平状态）；\\
  沿 \(y=x\) 对折后，反函数 \(x=y^{1/3}\) 在 \(y=0\) 处被``拉竖''（切线呈竖直状态）。由于竖直切线的斜率趋于无穷大，故反函数在原点不可导。
\end{itemize}

\subsection{小结：局部线性化的统一视角}

复合求导、隐式求导与反函数求导，本质上是\textbf{链式法则在不同几何约束与变量依赖关系下的统一体现}：
\begin{itemize}\tightlist
  \item \textbf{复合求导}：变量呈单向链条关系，导数沿着信息传递链逐层连乘；
  \item \textbf{隐式求导}：变量间被约束方程绑定，通过对共同自变量求导维持切向运动；
  \item \textbf{反函数求导}：自变量与因变量互换角色，导数转化为对应工作点上局部放大率的倒数。
\end{itemize}

上述法则均属于\textbf{局部（Local）}分析——它们精确描述了函数在单一工作点附近的“微观线性放大”行为。要了解这些局部伸缩如何在整个区间上累积并决定宏观变化量，则需要引入下一节的核心工具：中值定理与微积分基本定理。